\section{Four levels that must not be conflated}
\label{sec:levels}

The word ``representation'' can connect questions about reality, meaning,
warrant, and proof.  Moving among these levels without notice makes an
apparently deep thesis either trivial or overstated.  We therefore separate
four claims.

\subsection{Ontology}

An ontological claim concerns what exists or is fundamental: objects,
relations, events, fields, algebraic structures, or something else.  The
formalism developed below does not select among these options.  It supports a
more modest thought: when one physical theory is realized through several
descriptions, the content for which cross-description recovery can be proved
is a particularly robust candidate for objective commitment.  That is
compatible with ontic or epistemic structural realism, entity realism,
constructive empiricism, and pragmatism.  It does not imply that the universe
literally is a diagram or that relations exist without relata.

\subsection{Semantics}

A semantic claim specifies what mathematical items represent and under which
interpretation.  A gauge potential, an orbit, a Wilson loop, a probability
measure, a blocked action, a Schwinger functional, and a spectral projection
do not represent the same type of thing.  Their physical content depends on
an observable algebra, states that assign expectation values, a redundancy
relation, and an interpretation of dynamics and units.  Two structures can
be isomorphic while receiving different interpretations; conversely,
nonisomorphic models can agree on a restricted empirical domain
\cite{Coffey2014,DeHaroButterfield2018}.  The calculus consequently types both
formal maps and the claims transported by them.

\subsection{Epistemology}

The epistemic issue is when a result in one description warrants a result in
another.  The answer is property-relative.  A gauge quotient can be exact for
invariant expectation values while losing manifest local sections.  A coarse
graining can approximate long-distance observables while erasing a boundary
state needed for conditional estimates.  A Euclidean functional can encode a
Hamiltonian spectrum only when positivity and reconstruction hypotheses hold.
There is no contradiction in assigning different fidelity verdicts to the
same arrow, because the predicates have different bearers and inferential
roles.

\subsection{Proof methodology}

The methodological claim is operational.  A proof that crosses a
representation interface must record:
\begin{enumerate}[label=(\roman*),itemsep=2pt]
  \item the source and target objects;
  \item the observable or state map and its domain;
  \item the target property and all of its prerequisites;
  \item exact equations or approximation moduli;
  \item structures unavailable in the target variables;
  \item separate recovery theorems; and
  \item the order and uniformity of all limiting quantifiers.
\end{enumerate}
This is the level on which the paper's formal results operate.  The utility of
the calculus does not wait on settlement of the ontology.

\subsection{A dispensable use of information}

The phrase ``physical information'' is often asked to perform work belonging
to all four levels.  We use it only as shorthand for distinctions among
states that are recoverable by a declared family of physical observables.  Its
bearer is a state or correlation functional; its equivalence relation is
agreement on the physical algebra; its accessibility is the state--observable
pairing; and its dynamics is a typed action on those pairings.  Gauge change,
physical time evolution, RG transformation, and OS reconstruction are thus
different operations.  If bearer, equivalence relation, and inferential role
are absent, the word ``information'' can be deleted without changing the
framework.  This prevents an informational slogan from substituting for a
physical ontology.
